\section{Introduction}
\label{sec:introduction}

Classical models of romantic selection and matching assume a total order over
prospective partners: there is a first choice, a second choice, a third choice,
and so on~\citep{gale1962college,roth1990two}. Once an agent establishes that
total order, commitment is modeled as matching with the highest-ranked available
candidate, while all remaining candidates sit cleanly below on a single ordinal
ladder. Such models are mathematically convenient, intuitive, and, we argue,
empirically inadequate for a widespread class of contemporary relational
configurations.

Consider an individual $W$ who maintains regular contact with several people
$G_1, G_2, \dots, G_n$ at the same time. The connections are not interchangeable
substitutes waiting in a queue. Instead, they are functionally differentiated.
One partner provides public visibility and shared weekend social appearances.
Another provides late-night emotional processing and receives long voice notes
at 2~a.m. A third provides practical coordination, meals, or financial support.
A fourth is maintained at low cadence as an active contingency. If an external
observer, or one of the partners, asks $W$ a seemingly elementary question, namely
which of $\{G_1, \dots, G_n\}$ is \emph{the one}, the observed response is
characteristically evasive, delayed, or deflected with a counter-question.

Standard preference theory treats this evasiveness as measurement error,
incomplete introspection, or strategic dishonesty. \emph{Relational Superposition
Theory} (\RST) takes the opposite starting point: the evasiveness is not noise
obscuring a hidden linear ranking. It is the observable signature of a system in
which no linear ranking exists prior to inquiry. Because each partner occupies a
distinct functional niche and holds a nontrivial share of $W$'s relational
investment, the question ``who is first?'' has no well-defined answer until an
intervention forces the system to select one.

\subsection{Why Functional Partitioning Breaks Ordinal Rank}
\label{subsec:intro-partition}

In a single-attribute decision problem, ordering candidates requires only a scalar
comparison. In a multi-partner relational system, however, value is distributed
across incommensurable functional dimensions: public legitimacy, emotional
attunement, material reliability, novelty, and availability. Optimizing one
dimension frequently conflicts with optimizing another. Rather than collapsing
those competing criteria onto a single person before commitment, a focal agent
with sufficient options can partition functions across several concurrent
connections.

Once partitioning occurs, two structural consequences follow immediately. First,
no single partner achieves unit primary weight; each partner holds an intermediate
weight $w_i \in (0, 1)$ reflecting his share of attention and priority. Second,
the question ``where do I stand?'' (rank) and the question ``what am I to you?''
(role) cease to be jointly answerable with arbitrary precision. Pinning a partner
firmly to a specialized role such as emotional confidant leaves his ordinal rank
against the weekend companion or financial provider undefined. Conversely, forcing
a strict rank declaration smears the chosen partner's role across obligations he
did not originally occupy.

\subsection{The Role of Perturbation and Social Coupling}
\label{subsec:intro-perturbation}

A superposed relational system can persist smoothly for weeks or months while
routine communication continues. Its stability depends on keeping rank queries
unasked and keeping partners separated in the broader social graph. Two distinct
mechanisms break that stability:

\begin{enumerate}[nosep]
  \item \textbf{Internal perturbation (the poke).} Any direct status inquiry,
        calendar confrontation, or boundary demand forces $W$ to resolve the
        ambiguity in favor of one partner, collapsing the prior superposition
        probabilistically according to pre-poke weights.
  \item \textbf{External coupling (decoherence cascade).} If partners discover
        one another through shared friends, group chats, or overlapping social
        venues and compare notes, collapse propagates across the friend graph.
        Rather than leaving one partner in the primary state, all active weights
        can decay to zero simultaneously.
\end{enumerate}

The persistence of superposition is also unevenly distributed across populations.
Maintaining several concurrent connections requires a supply of viable partners
who tolerate partial availability. That supply varies systematically with
perceived attractiveness and social reach, motivating the structural link between
optionality and superposition rather than purely moral or personality-based
explanations.

\subsection{Summary of Contributions}
\label{subsec:intro-contributions}

This monograph develops \RST{} as a self-contained mathematical and computational
framework. Its specific contributions are:
\begin{enumerate}[nosep]
  \item \textbf{Order-theoretic diagnosis (\cref{sec:background}).} We show how
        functional role partitioning induces strict comparability failures that
        prevent representation by a single weak or total order.
  \item \textbf{State-vector formalism (\cref{sec:framework}).} We construct the
        finite-dimensional partner state space $\vspaceR$ and system state vector
        $\ket{\Psi}$, define relational superposition, and prove that no partner
        occupies rank one prior to collapse.
  \item \textbf{Operator derivation of uncertainty (\cref{sec:uncertainty}).} We
        formulate rank and role as non-commuting Hermitian operators $\opR$ and
        $\opL$, derive the Rank-Role Inequality $\Delta R \cdot \Delta L \ge \kappaS$
        via Cauchy-Schwarz, characterize minimum-uncertainty states, and work out
        exact $2\times 2$ and $4\times 4$ examples.
  \item \textbf{Attention-drift dynamics (\cref{sec:dynamics}).} We model weekday,
        weekend, and emotional-stress transitions as a time-inhomogeneous Markov
        chain and characterize its stationary distribution.
  \item \textbf{Poke collapse and social percolation (\cref{sec:poke,sec:cascade}).}
        We formalize the Poke Rule, post-collapse renormalization, irreversibility,
        and a bond-percolation threshold $q_c$ governing social-graph cascades.
  \item \textbf{Optionality and symmetry (\cref{sec:conjecture,sec:symmetry}).}
        We formalize the Attractiveness-Optionality Conjecture under a binomial
        retention model, prove monotonicity, propose an unadministered empirical
        design, and establish gender symmetry.
  \item \textbf{Reproducible simulation and formal proofs (\cref{sec:simulations,sec:verification}).}
        We report deterministic synthetic Monte Carlo studies across $2{,}400$ operator
        samples, $90$-day attention trajectories, optionality curves, and graph
        cascades, alongside machine-checked Lean~4 proofs of normalization,
        absence of a primary winner, optionality monotonicity, and the finite-dimensional
        uncertainty bound.
\end{enumerate}

\begin{remark}[Pronoun convention]
\label{rem:pronouns}
Following the initial observational motivating case, we refer to the focal person
as $W$ and to members of the partner set $\{G_i\}$ using the male pronoun for
notational consistency. \cref{sec:symmetry} proves that every equation, bound,
and dynamical transition in the theory is invariant under relabeling of partners
and roles, including the mirror configuration in which a male focal agent
maintains multiple female partners.
\end{remark}
